% =====================================================================
% Agentic FedGrid -- Abstract, Contributions, Results, Conclusion.
% Every number below is from the EC2 runs in paper_figures/aws_logs/
% (parsed by build_results.py into results_tables.tex / results_final.json).
% Tables referenced: tab:backbone, tab:refs, tab:fed_rounds_v1, tab:fed_rounds_v2,
% tab:operating, tab:threshold, tab:ablation, tab:policy_ton, tab:policy_unsw,
% tab:agent_ton, tab:agent_unsw, tab:analyst.
% Figures: fig1_infrastructure.png, fig2_rounds.png, fig3_ablation_bars.png,
% fig4_agent_bars.png.
% =====================================================================

\begin{abstract}
Network intrusion detection at the edge must adapt to attacks that only some
sites have seen, without moving raw traffic off the device and without letting
one bad update or one false alarm disrupt legitimate hosts. We present
Agentic FedGrid, a deployed system in which autonomous agents on edge nodes
detect, respond to, and learn from network flows. A frozen tabular Transformer
is adapted with low-rank (LoRA) adapters federated through a serverless AWS
control plane (10.3\,KB per upload vs.\ 103.7\,KB for the full model). The
agent chooses actions by expected cost, escalates blocks gradually with
automatic expiry, learns from analyst rollbacks, shares and retracts threat
intelligence with its peers, and adopts a new global model only if it does
not degrade its own held-out performance. All experiments ran on four EC2
nodes with multiple seeds.

Our evaluation is deliberately threshold-free where it matters. Four
federated rounds raise recall on an attack class withheld from pretraining
(Fuzzers) from 41.1\% to 54.5\% at the default threshold, but the
false-positive rate rises from 14.3\% to 19.2\%. Moving only the decision
threshold of the frozen backbone to the same FPR yields 53.4\%, and the
Fuzzers-vs-benign AUROC does not improve (0.823 vs.\ 0.818). The apparent
gain is an operating-point shift, a pitfall we expect to affect other
federated IDS evaluations. Spreading the class over three clients or running
5$\times$ more local steps does not change this: recall and FPR rise together
(up to 87.2\% at 37.4\% FPR), the Fuzzers AUROC stays at 0.817--0.820, and the
frozen model with a moved threshold matches or beats every variant.
Full-model federation degrades the ranking (AUROC 0.962 to 0.950), while LoRA
at the default local budget preserves it (0.959). The robustness
results hold: under a model-replacement attack FedAvg collapses to 47.5\%
accuracy, whereas coordinate-wise median aggregation holds 86.5\% and the
adoption gate rejects every poisoned model. At the decision layer, the full
agent cuts benign traffic dropped by blocking from 49.8\% to 12.5\% (ToN\_IoT)
and from 22.1\% to 2.1\% (UNSW-NB15) while letting fewer attacks through. With
an imperfect analyst (10\% wrong verdicts) mistaken rollbacks let attacks
through; shortening the learned exception from 30 to 5 minutes cuts that
leakage by about two-thirds. The backbone reaches 87.7\% accuracy on the official
UNSW-NB15 test set (gradient-boosted trees: 90.1\%), 99.3\% on ToN\_IoT, and
1.7\,ms inference on the edge node.
\end{abstract}

\section*{Contributions}
\begin{enumerate}
\item A deployed edge--cloud architecture that federates LoRA adapters of a
tabular Transformer IDS through serverless AWS services (API Gateway, Lambda,
DynamoDB, S3), with straggler handling and lock-based round coordination. Raw
flows never leave the device. The live backend closed 12{,}610 rounds over
25 days.
\item An edge agent whose actions go beyond classification: cost-aware action
selection, graded and self-expiring blocks, analyst-feedback rollback with
learned exceptions, peer threat-intel sharing with retraction, and verified
adoption of global models. It is evaluated against both an ideal and an
imperfect analyst.
\item A threshold-free evaluation of federated adaptation showing that, in
this setting, the recall gain on an unseen class is fully explained by a
shift of the operating point, including when the class is spread over three
clients or local training is 5$\times$ longer. This is a reporting pitfall for
federated IDS work that uses default-threshold recall.
\item A multi-seed evaluation on real cloud hardware: LoRA rank vs.\
full-model federation, FedAvg vs.\ FedProx vs.\ median aggregation,
hard-example mining, adversarial training, model poisoning, and the agent's
decision layer. Gradient-boosted and centralized references are included.
\end{enumerate}

\section{Results}

\subsection{Backbone and reference models}
Table~\ref{tab:backbone} reports the frozen backbone before any federation.
Moving from one token for all continuous features (v1) to one token per
feature with a longer one-cycle schedule (v2) raises UNSW-NB15 accuracy from
87.19\% to 87.73$\pm$0.35\%, lowers its FPR from 14.31\% to
12.94$\pm$1.13\%, and lifts ToN\_IoT to 99.28$\pm$0.04\% (FPR 1.64\%).
Inference stays at 1.7\,ms per flow on a t3.small. The better general fit
comes with lower recall on the withheld Fuzzers class before federation
(36.81\% vs.\ 41.06\%), so v2 is not uniformly better.

Table~\ref{tab:refs} places the backbone between two references trained on
the same eight features. Gradient-boosted trees trained on the same
pretraining pool reach 90.10$\pm$0.17\% accuracy at 8.9\% FPR, so the
Transformer itself costs about 2.4 points on this split. We keep it because
LoRA federation, FGSM training and attention attribution all require a
differentiable model, and the tree's advantage does not extend to the unseen
class (34.3\% Fuzzers recall vs.\ 41.1\%). The backbone scores 97.9\% on a
random split of its own pretraining pool. The roughly 10-point drop to the
official test file is distribution shift between the official training and
test files, not a lack of fit. A Transformer trained
centrally on all device data, Fuzzers included, reaches 89.8\% Fuzzers recall
but only 82.4\% accuracy at 36.4\% FPR (3 seeds).

\subsection{Federated adaptation to an unseen attack class}
Fuzzers flows are excluded from pretraining and appear only on device~1,
revealed 25/50/75/100\% over four rounds. All metrics are measured on the
global model on device~2, which never sees a Fuzzers flow.

\textbf{At the default threshold} (Table~\ref{tab:fed_rounds_v1},
Fig.~\ref{fig:rounds}), Fuzzers recall rises from 41.06\% to
54.50$\pm$1.92\% (backbone v1, 4 seeds) and replay-set accuracy from 28.14\% to
43.53\%. Accuracy moves only from 87.19\% to 86.46\%, but the FPR rises from
14.31\% to 19.21\%. Backbone v2 gives 36.81\%~$\rightarrow$~47.39$\pm$9.83\%
at 17.01\% FPR (Table~\ref{tab:fed_rounds_v2}).

\textbf{Threshold-free} (Tables~\ref{tab:operating},~\ref{tab:threshold}), the
picture changes. Across rounds, AUROC stays at 0.962~$\rightarrow$~0.959, the
Fuzzers-vs-benign AUROC at 0.823~$\rightarrow$~0.818, and Fuzzers recall at a
fixed 10\% FPR falls from 38.0\% to 34.8\%. Moving only the frozen backbone's
threshold to the FPR each federated variant ends at gives the same recall:
53.4\% vs.\ 54.5\% at 19.2\% FPR, 51.6\% vs.\ 50.5\% at 18.3\%, and 57.0\% vs.\
59.1\% at 21.2\%. Federation in this setting therefore moves the operating
point toward ``attack''; it does not teach the model to separate Fuzzers from
benign traffic. The information is not absent. The centrally trained model
reaches 0.847 Fuzzers AUROC and 48.1\% recall at 10\% FPR (one seed).

\textbf{More clients and longer local training.} Two obvious explanations
are that only one of four clients holds the class and that each round runs
only 100 local steps. We tested both on the same Fuzzers flows: split into
disjoint shards over two or three clients, with 5$\times$ more local steps
(500 per round), and both combined (Tables~\ref{tab:ablation},~\ref{tab:operating},~\ref{tab:threshold};
3 seeds each). Default-threshold recall rises to 51.1\% (two clients),
71.4\% (5$\times$ steps), 79.4\% (three clients) and 87.2\% (both), but the
FPR rises with it, to 18.5\%, 28.2\%, 33.6\% and 37.4\%. The Fuzzers-vs-benign
AUROC stays at 0.817--0.820 in every variant (0.823 before federation), and
recall at 10\% FPR stays at 33.8--36.1\% (38.0\% before). The frozen backbone
with only its threshold moved to the same FPR matches or beats each variant:
51.9\% vs.\ 51.1\%, 71.7\% vs.\ 71.4\%, 85.1\% vs.\ 79.4\%, and 89.7\% vs.\
87.2\%. Longer local training also costs robustness (FGSM accuracy 74.5\%
$\rightarrow$ 53.9\%), and with three clients and 5$\times$ steps two of three
seeds collapse (UNSW accuracy 59.8\% and 59.6\%, AUROC 0.77 and 0.75).
Participation and local budget are therefore not the bottleneck in the regimes
we tested. Together with the full-model result below, which removes the
capacity limit and degrades the ranking further, this suggests that the
binary local objective itself moves the decision boundary rather than the
representation. We regard default-threshold recall as an unreliable measure
of federated gains and report both.

\subsection{Ablations}
Table~\ref{tab:ablation} and Fig.~\ref{fig:ablation} vary one component at a
time. Unless stated otherwise, each variant uses backbone v1, LoRA $r{=}4$ and
FedAvg, with 3 seeds.

\textbf{Rank and full-model federation.} Ranks 1, 2 and 4 behave almost
identically: Fuzzers recall 50.5--52.4\%, FPR $\approx$18.3\%, at 2.8--10.3\,KB
per upload. Rank 8 (20.3\,KB) reaches 59.1\% at 21.2\% FPR. Full-model FedAvg
sends 10$\times$ more data (103.7\,KB) and reaches 78.0\% recall at 33.1\% FPR.
Threshold-free, the full model is worse: AUROC 0.950 vs.\ 0.959 for LoRA,
30.2\% vs.\ 34.8\% Fuzzers recall at 10\% FPR, and 78.0\% vs.\ 84.0\% for the
frozen backbone at the same 33.1\% FPR. Updating every weight from one
site's skewed data degrades the shared ranking. The low-rank restriction
limits this degradation, but it does not produce a gain either.

\textbf{Local objective.} FedProx ($\mu{=}0.01$) gives no measurable gain over
FedAvg: 52.4\% vs.\ 50.5\% recall and 19.0\% vs.\ 18.3\% FPR, both within
one standard deviation. This is consistent with the already bounded local
update, which limits client drift without a proximal term.

\textbf{Training components.} Removing FGSM training leaves clean accuracy
essentially unchanged (86.78\% vs.\ 86.47\%). Accuracy under an FGSM
perturbation ($\epsilon{=}0.05$) falls from 74.49\% to 52.37\%. Removing
hard-example mining does not change Fuzzers recall (52.97\%), but robustness
becomes unstable across seeds (58.7$\pm$12.8\%). We report mining as having
no clear benefit.

\textbf{Poisoning.} One client uploads $-10\times$ its honest adapter
(model replacement). Under FedAvg the global model collapses to near-chance
accuracy (47.5$\pm$13.1\%), alternating between labelling everything benign
and everything malicious. With the adoption gate, every device rejected every
poisoned global model in all rounds and seeds. However, the isolated devices
keep training only on their own skewed data, and device~2's FPR drifts to
38.8\%. Coordinate-wise median aggregation removes the attack at the server:
86.54$\pm$0.21\% accuracy and 18.46$\pm$2.20\% FPR, matching the clean runs.
Median aggregation on clean data costs little (86.31\% accuracy, 20.1\% FPR;
AUROC 0.959, identical to FedAvg). The two defences are complementary: the
median protects the shared model, and the gate protects a device if the
server is compromised.

\subsection{The agentic decision layer}
\textbf{Action selection} (Tables~\ref{tab:policy_unsw},~\ref{tab:policy_ton}).
On UNSW-NB15, the cost-aware rule cuts missed attacks from 11.6\% to 2.4\% of
attack flows and lowers the average cost per flow by 23\%. The price is that
it alerts on 31.5\% of benign flows instead of 12.6\%, which means more
analyst workload. On ToN\_IoT both rules are already near-perfect, and the
cost rule halves missed attacks (0.92\%~$\rightarrow$~0.39\%).

\textbf{Response and learning} (Tables~\ref{tab:agent_ton},~\ref{tab:agent_unsw};
Fig.~\ref{fig:agent}). Blocking per source IP magnifies false positives: one
wrong verdict on a busy benign host cuts off all of its later traffic. With
permanent blocks (configuration A) this drops 49.8\% of benign ToN\_IoT flows
and 22.1\% of benign UNSW-NB15 flows.
\begin{itemize}
\item Graded, self-expiring blocks (C) reduce this to 28.6\% and 1.4\%,
without any analyst.
\item Analyst rollback with learned exceptions (D) reduces it further, to 9.4\% and 1.2\%.
\item Peer intelligence with retraction (E) costs a little benign traffic
(12.5\% and 2.1\%) but stops more attacks. Attack flows that get through fall
from 0.029\% to 0.019\% (ToN\_IoT) and from 0.38\% to 0.13\% (UNSW-NB15)
relative to the threshold-and-block baseline.
\end{itemize}
With an oracle analyst, the full agent therefore dominates the baseline on
both axes and both datasets.

\textbf{Imperfect analyst} (Table~\ref{tab:analyst}). In the next experiment,
the analyst reviews only a random half of the blocks and is wrong 10\% or
20\% of the time. Benign traffic dropped stays well below the baseline
(24.0\% vs.\ 49.8\% on ToN\_IoT, 4.1\% vs.\ 22.1\% on UNSW-NB15 at 10\% error).
Attack leakage, however, grows sharply: 4.5\% of ToN\_IoT attack flows get
through at 10\% error and 11.8\% at 20\% error, against 0.083\% for
configuration C with no analyst at all. The cause is the learned exception. A
single wrong rollback of an attacker grants it 30 minutes in which block
verdicts are downgraded to alerts, and the retraction propagates to every
peer. Shortening the exception to 5 minutes cuts the leakage by about
two-thirds (4.5\%~$\rightarrow$~1.5\% on ToN\_IoT and 3.1\%~$\rightarrow$~1.1\%
on UNSW-NB15 at 10\% error). It costs more benign blocking with an ideal
analyst (12.5\%~$\rightarrow$~18.6\% on ToN\_IoT). Analyst feedback is only
worth enabling where analyst verdicts are reliable. Otherwise graded blocks
alone (C) are the safer configuration.

\section{Conclusion}
Agentic FedGrid shows that low-rank federated adaptation and an agent that
reasons about the cost and trustworthiness of its own actions can be combined
on commodity cloud hardware. The measurements support four conclusions:
\begin{itemize}
\item Default-threshold recall overstates federated adaptation. In our
setting, the recall gain on an unseen class is matched by moving the frozen
model's threshold, and the class separation (AUROC) does not improve. This
held with the class spread over up to three clients and with 5$\times$ longer
local training. Federated IDS results should be reported threshold-free.
\item Low-rank adaptation does not add discrimination here, but at the default
local budget it preserves the backbone's ranking where full-model federation
degrades it. Median
aggregation together with verified adoption neutralises model-replacement
poisoning.
\item The main obstacle in edge intrusion detection is false positives, and
per-host blocking magnifies them. Graded, self-expiring response removes most
of this cost without any human input.
\item Human feedback helps only when it is reliable. Mistaken rollbacks are
amplified by learned exceptions and shared retractions, so their duration
must be short.
\end{itemize}
Future work: federated objectives that change the representation rather than
the threshold, such as contrastive or multi-class local objectives (more
clients holding the class and longer local training did not help here); real
analyst feedback; and temporal (multi-flow)
models that reduce the per-flow false positives behind most of the costs
measured here.
